% Drawing 156-4/03-4 -- folder 156-4, batch 03, archive page 51 (PDF page 52).
% Transcription: moitié gauche, troisième schéma : $\overline{F}' \subset \widetilde{F} \xrightarrow{\sim} \vec{F} \xrightarrow{\sim} F$ au-dessus de $\widetilde{G} \xrightarrow{\sim} \overline{G} \xrightarrow{\sim} G$, avec des flèches verticales.
% Kind: relation diagram (subdivisions ~-> and inclusions). Method 2 (schematic redraw from the facsimile). Status: partial.
% Notes: Top row: F-bar' (included in) F-tilde ~-> F-vec ~-> F; bottom row: G-tilde ~-> G-bar ~-> G, each node under the one above it, joined by a small vertical cup under F-tilde and under F-vec and a short descending arrow under F. Labels after the transcription; doubtful: the middle bottom node (read G-bar; on the crop it may carry a tilde as well as a bar). Left out: the top of the large arc that encloses the fourth schema, which passes under this one. Nothing struck. Method: a relation diagram of formulas, which the transcription skill would set as text (tikz-cd); since it stands only as a \drawing mark, it is set here as a TikZ layout of the formula: every symbol is a LaTeX text node at its place on the page, the arrows are drawn between them. Redrawn 2026-11-09 (issue 128): to check again. Question settled by a machine pass on the facsimile, 2026-11-12 (maths): under F-tilde stand a small cup, a tilde and G, and under F-vec a small cup, a bar and G: the mark above the bar is the vertical inclusion cup, not a tilde, and the cube below has the same row G-tilde ~-> G-bar ~-> G; the middle bottom node is G-bar, as drawn. Not checked by a reader.
% Coordinates: 1 unit = 1 mm = 1/100 of the crop width (bbox in page fractions: [0.2, 0.645, 0.48, 0.72]).
% Redrawn by Claude (claude-opus-5-5), issue #73, 2026-11-09; not checked by a human; not compiled (no TeX in the session).
% Rights: reproduces part of the work -- behind RENDER_TRANSCRIPTS, never published by itself.
\documentclass[tikz,border=3pt]{standalone}
\usepackage{amsmath,amssymb}
\usetikzlibrary{patterns,arrows.meta}
\tikzset{pencil/.style={black!45}, struck/.style={black!35}, mark/.style={circle,draw,inner sep=1.3pt}, pt/.style={circle,fill,inner sep=0.9pt}, opt/.style={circle,draw,fill=white,inner sep=1.1pt}}
\begin{document}
\begin{tikzpicture}[x=1mm,y=1mm,line cap=round,line join=round,font=\scriptsize]
\node[inner sep=1pt] at (34,-15) {\small $\overline{F}'$};
\node[inner sep=1pt] at (45,-15) {\small $\subset$};
\node[inner sep=1pt] at (56,-14) {\small $\widetilde{F}$};
\draw[-{Stealth[length=2mm]}] (61,-14) -- (70.5,-14);
\node[inner sep=1pt] at (65.8,-11.4) {\small $\sim$};
\node[inner sep=1pt] at (75.5,-14) {\small $\vec{F}$};
\draw[-{Stealth[length=2mm]}] (80.5,-13.5) -- (88,-13.5);
\node[inner sep=1pt] at (84.2,-10.9) {\small $\sim$};
\node[inner sep=1pt] at (93,-13) {\small $F$};
\node[inner sep=1pt] at (56,-21) {\small $\cup$};
\node[inner sep=1pt] at (75.5,-19.5) {\small $\cup$};
\draw[-{Stealth[length=2mm]}] (93.5,-17) -- (93.5,-21);
\node[inner sep=1pt] at (56,-30) {\small $\widetilde{G}$};
\draw[-{Stealth[length=2mm]}] (61.5,-29.5) -- (71,-28);
\node[inner sep=1pt] at (66.2,-26.1) {\small $\sim$};
\node[inner sep=1pt] at (75.5,-26.5) {\small $\overline{G}$};
\draw[-{Stealth[length=2mm]}] (80.5,-26) -- (88,-25.5);
\node[inner sep=1pt] at (84.2,-23.1) {\small $\sim$};
\node[inner sep=1pt] at (95,-24) {\small $G$};
\end{tikzpicture}
\end{document}
